\documentclass[10pt,a4paper]{amsart}
\usepackage[margin=19mm]{geometry}
\usepackage{amsmath,amssymb,mathtools}
\usepackage[scr=rsfs,cal=euler]{mathalfa}
\usepackage{xfrac}
\usepackage[unicode,hidelinks,bookmarks=false]{hyperref}
\newcommand{\ie}{i.e.\ }
\newcommand{\cf}{cf.\ }
\newcommand{\resp}{respectively\ }
\newcommand{\Subsection}{\subsection}
\providecommand{\nobreakdash}{\nobreak\hbox{-}}
\setlength{\emergencystretch}{2em}
\multlinegap0pt
\usepackage{amssymb}
\usepackage{float}
\usepackage{csquotes}
\usepackage{tikz}
\newtheorem{proposition}{Proposition}
\title{\Large\bfseries Linearization Problem for a System of \\Two Second-Order ODEs via Cartan's \\Method: Branch I}
\author{Batoul M. Raddad, Ahmad Y. Al-Dweik, Marwan Aloqeili, F. M. Mahomed}
\date{Source: arXiv:2607.16840v1 (CC BY 4.0)}
\begin{document}
\maketitle

\noindent\emph{Source and licence.} \Large\bfseries Linearization Problem for a System of \\Two Second-Order ODEs via Cartan's \\Method: Branch I. Version arXiv:2607.16840v1 (CC BY 4.0), distributed by arXiv under the Creative Commons Attribution 4.0 International licence, http://creativecommons.org/licenses/by/4.0/, as recorded in the arXiv licence metadata for this exact version and shown on the abstract page https://arxiv.org/abs/2607.16840v1. Full licence notice: \texttt{licenses/arxiv-2607.16840-CC-BY-4.0.txt}.\par\smallskip

\noindent\textit{Editorial scope.} Counted unit: the article's proposition s3-8 (source0.tex lines 1054--1105). The article's complete original proof of it is delivered with it (source0.tex lines 1106--1229, one \texttt{proof} environment of 124 lines). No mathematics is written, repaired or re-derived: the statement and the proof are the article's own lines, in the article's own notation, and no retained line is substituted or re-flowed.  Retained uncounted as the source-local context the proof needs: lines 383--392; lines 1007--1020; lines 1022--1042.  Nothing was omitted from any retained span, so the package carries no omission record.  No retained line contains a citation, so the wrapper carries no bibliography and no bibliography entry.  No retained line contains a figure environment, an image-inclusion command or drawing code, so no image is delivered and the package declares no asset dependency.  Editorial formatting only, in addition: an AMS \texttt{amsart} preamble that restates the article's own declarations and macros used by the retained lines, namely the environments \texttt{proposition} on separate counters, together with the standard packages those lines need.  The retained environments print in this document as Proposition 1 in order of appearance, whereas the article prints the same items with the numbering of its own full document.  The complete native source of the article as distributed in the arXiv e-print for this exact version is bundled unchanged as \texttt{sources/205505/source0.tex}.
\begin{equation}\label{87}
\Omega=
\begin{pmatrix}
1&0&0&0&0\\
0&1&0&0&0\\
I_1&I_2&1&0&0\\
I_3&I_4&0&1&0\\
0&0&0&0&1
\end{pmatrix},
\end{equation}

\begin{equation}\label{coframe}
\begin{pmatrix}
\theta\\
\alpha\\
\end{pmatrix}=
\begin{pmatrix}
S&0\\
\Gamma&\Lambda\\
\end{pmatrix}
\begin{pmatrix}
\omega\\
V
\end{pmatrix},
\end{equation}

\begin{equation}\label{s3-10}
\begin{aligned}
&S=\begin{pmatrix}
a_1&\frac{J_2a_1}{J_3}&0&0&0\\
a_3&\frac{J_2a_3+a_1a_{15}^2}{J_3}&0&0&0\\
\frac{K_4a_1}{2a_{15}}&\frac{J_2K_4a_1}{2J_3a_{15}}&\frac{a_1}{a_{15}}&\frac{J_2a_1}{J_3a_{15}}&0\\
\frac{K_4a_3}{2a_{15}}&\frac{(J_2a_3+a_1a_{15}^2)K_4}{2J_3a_{15}}&\frac{a_3}{a_{15}}&\frac{J_2a_3+a_1a_{15}^2}{J_3a_{15}}&0\\
-2K_7a_{15}&\frac{2(K_5-J_2K_7)a_{15}}{J_3}&0&0&a_{15}
\end{pmatrix},\\
&\Lambda=\begin{pmatrix}
\frac{1}{a_1}&0&\frac{2}{a_{15}}\\
-\frac{a_3}{a_1^2}&\frac{1}{a_1}&-\frac{2a_3}{a_1a_{15}}\\
0&0&\frac{1}{a_{15}}\\
\end{pmatrix},\\
 &\Gamma=\begin{pmatrix}
-\frac{J_2I_{3p_1}+J_3R_1}{J_3}&-\frac{J_2R_3+J_3R_5+2K_4K_5}{J_3}&2K_7&\frac{2(J_2K_7-K_5)}{J_3}&-\frac{2I_3J_2+J_3R_7}{2J_3}\\
\frac{2J_3(R_3-I_{3p_2})a_3-I_{3p_1}a_1a_{15}^2}{J_3a_1}&\frac{(2J_3R_4-2R_6)a_3-R_2a_1a_{15}^2}{J_3a_1}&-\frac{4K_7a_3}{a_1}&-\frac{2K_7a_1a_{15}^2+4R_6a_3}{J_3a_1}&\frac{2J_3K_4a_3-I_3a_1a_{15}^2}{J_3a_1}\\
I_{3p_2}+R_{2}&\frac{-J_3R_4+K_4R_6}{J_3}&2K_7&\frac{2R_6}{J_3}&-K_4
\end{pmatrix},
\end{aligned}
\end{equation}

\begin{proposition}\label{s3-8}
Assume that the system of two second-order ODEs
\begin{equation}\label{136}
u_1^{\prime\prime}=f^1(x,u_1,u_2,p_1,p_2),~~u_2^{\prime\prime}=f^2(x,u_1,u_2,p_1,p_2)
\end{equation}
 whose generalized Wilczynski invariant matrix $W$ has nonzero entries, is equivalent to the canonical form
\begin{equation}\label{138}
\bar{u}_1^{\prime\prime}=\bar{u}_1+\bar{u}_2,~~\bar{u}_2^{\prime\prime}=-(\bar{u}_1+\bar{u}_2)
\end{equation}
under the invertible point transformation
\begin{equation}\label{137}
\Phi:~~\bar{x}=\xi(x,u_j)~~\bar{u}_i=\phi_{i}(x,u_j).
\end{equation}
Then the linearizing point transformation (\ref{137}) where $\phi_{1u_1}-\xi_{u_1}\frac{D_x\phi_1}{D_x\xi}\neq0$
can be obtained by the following sequential procedure:\\
\textbf{Step 1.}  Determine the matrices 
\begin{equation}\label{s3-12}
B=(\bar{S})^{-1}S~\Omega, ~H=(\Lambda)^{-1}(\Gamma\Omega-\bar{\Gamma}B).
\end{equation}
 \textbf{Step 2.} Find the nonzero auxiliary functions $a_1,a_{15}$ and the auxiliary function $a_3$ by solving the system \\
\begin{equation}\label{168}
\begin{aligned}
&\begin{pmatrix}
a_{1u_1}&a_{1u_2}&a_{1p_1}&a_{1p_2}&D_xa_1\\
a_{3u_1}&a_{3u_2}&a_{3p_1}&a_{3p_2}&D_xa_3\\
a_{15u_1}&a_{15u_2}&a_{15p_1}&a_{15p_2}&D_xa_{15}\\
\end{pmatrix}=-H,
\end{aligned}
\end{equation}
\textbf{Step 3.} Find $\xi, \phi_i, \chi_i$ by solving the system
\begin{equation}\label{169}
\begin{aligned}
&\begin{pmatrix}
\phi_{iu_j}&\textbf{0}&D_x\phi_i\\
\chi_{iu_j}&\chi_{ip_j}&D_x\chi_i\\
\xi_{u_j}&\textbf{0}&D_x\xi
\end{pmatrix}=\begin{pmatrix}
I&\textbf{0}&\chi_i\\
\textbf{0}&I&\bar{f}^i\\
\textbf{0}&\textbf{0}&1\\
\end{pmatrix}B,
\end{aligned}
\end{equation}

where
\begin{equation}
\begin{aligned}
&\chi_i=\frac{D_x\phi_i}{D_x\xi},~~\bar{f}^1=\phi_1+\phi_2,~\bar{f}^2=-(\phi_1+\phi_2),\\
\end{aligned}
\end{equation}
 and $S,\Lambda,\Gamma$ are given in (\ref{s3-10}), $\Omega$ is given in (\ref{87}) and $\bar{S}, \bar{\Gamma}$ are the pullback of the matrices $S, \Gamma$ evaluated at the canonical form (\ref{138}), respectively, after substituting $\bar{a}_1=1, \bar{a}_3=0, \bar{a}_{15}=1$.
\end{proposition}

\begin{proof} 
Define the vectors $$U=(du_i-p_idx, ~dp_i-f^idx, ~dx)^{T},~~V=(da_1,da_3,da_{15})^{T}.$$ Then the symmetrical version of the coframe $(\theta,\alpha)$ given in (\ref{coframe}) provides
\begin{equation}
\begin{pmatrix}
\bar{S}&\textbf{0}\\
\bar{\Gamma}&\bar{\Lambda}\\
\end{pmatrix}\begin{pmatrix}
\bar{\Omega}&\textbf{0}\\
\textbf{0}&\text{I}\\
\end{pmatrix}\begin{pmatrix}
\bar{U}\\
\bar{V}\\
\end{pmatrix}=
\begin{pmatrix}
S&\textbf{0}\\
\Gamma&\Lambda\\
\end{pmatrix}\begin{pmatrix}
\Omega&\textbf{0}\\
\textbf{0}&\text{I}\\
\end{pmatrix}\begin{pmatrix}
U\\
V\\
\end{pmatrix},
\end{equation}
or equivalently,
\begin{equation}\label{141}
\begin{pmatrix}
\bar{U}\\
\bar{V}\\
\end{pmatrix}=
\begin{pmatrix}
(\bar{S}\bar{\Omega})^{-1}S\Omega&\textbf{0}\\
(\bar{\Lambda})^{-1}(\Gamma-\bar{\Gamma}(\bar{S})^{-1}S)\Omega&(\bar{\Lambda})^{-1}\Lambda\\
\end{pmatrix}\begin{pmatrix}
U\\
V\\
\end{pmatrix}.\\
\end{equation}
Substituting $\bar{a}_1=1, \bar{a}_3=0,\bar{a}_{15}=1$ in (\ref{141}) and then inserting the first prolongation of the point transformation (\ref{137}) with $\bar{p}_i=\chi_i=\frac{D_x\phi_i}{D_x\xi},$ results in
\begin{equation}\label{s3-11}
\begin{pmatrix}
\bar{U}\\
\textbf{0}\\
\end{pmatrix}=
\begin{pmatrix}
(\bar{S})^{-1}S\Omega&\textbf{0}\\
(\bar{\Lambda})^{-1}(\Gamma-\bar{\Gamma}(\bar{S})^{-1}S)\Omega&(\bar{\Lambda})^{-1}\Lambda\\
\end{pmatrix}\begin{pmatrix}
U\\
V\\
\end{pmatrix},
\end{equation}
where $\bar{U}=(d\phi_i-\chi_id\xi,~d\chi_i-\bar{f}^id\xi, ~d\xi)^{T}$ and
\begin{equation}
\bar{S}=\begin{pmatrix}
1&1&0&0&0\\
0&1&0&0&0\\
0&0&1&1&0\\
0&0&0&1&0\\
0&0&0&0&1
\end{pmatrix}, \bar{\Gamma}=\begin{pmatrix}
0&0&0&0&0\\
0&0&0&0&0\\
0&0&0&0&0\\
\end{pmatrix}, \bar{\Lambda}=\begin{pmatrix}
1&0&2\\
0&1&0\\
0&0&1
\end{pmatrix}.
\end{equation}
 Multiplying both sides of (\ref{s3-11}) by $\begin{pmatrix}
\mathbf{I}&0\\
0&\Lambda^{-1}\end{pmatrix}$  provides
\begin{equation}\label{142}
\begin{pmatrix}
\bar{U}\\
\textbf{0}\\
\end{pmatrix}=
\begin{pmatrix}
B&\textbf{0}\\
H&\textbf{I}\\
\end{pmatrix}\begin{pmatrix}
U\\
V\\
\end{pmatrix},
\end{equation}
where $B=(\bar{S})^{-1}S~\Omega, ~H=(\Lambda)^{-1}(\Gamma-\bar{\Gamma}(\bar{S})^{-1}S)\Omega$.\\

 Furthermore, using 
\begin{equation}
V=\begin{pmatrix}
a_{1u_1}&a_{1u_2}&a_{1p_1}&a_{1p_2}&D_xa_1\\
a_{3u_1}&a_{3u_2}&a_{3p_1}&a_{3p_2}&D_xa_3\\
a_{15u_1}&a_{15u_2}&a_{15p_1}&a_{15p_2}&D_xa_{15}\\
\end{pmatrix}U,
\end{equation}
together with (\ref{142}) provides (\ref{168}). 
Also, invoking
\begin{equation}
\overline{\vphantom{x}U}=\begin{pmatrix}
I&\textbf{0}&-\chi_i\\
\textbf{0}&I&-\bar{f}^i\\
\textbf{0}&\textbf{0}&1\\
\end{pmatrix}\begin{pmatrix}
d\phi_i\\
d\chi_i\\
d\xi
\end{pmatrix},
\end{equation}
one gets
\begin{equation}\label{167}
\overline{\vphantom{x}U}=\begin{pmatrix}
I&\textbf{0}&-\chi_i\\
\textbf{0}&I&-\bar{f}^i\\
\textbf{0}&\textbf{0}&1\\
\end{pmatrix}\begin{pmatrix}
\phi_{iu_j}&\textbf{0}&D_x\phi_i\\
\chi_{iu_j}&\chi_{ip_j}&D_x\chi_i\\
\xi_{u_j}&\textbf{0}&D_x\xi
\end{pmatrix}U.
\end{equation}
Hence, (\ref{142}) and (\ref{167}) provides (\ref{169}). 
This completes the proof.
\end{proof}

\end{document}
